\section{讨论}

路由实验回答第一个研究问题。DeepSeek Flash和GLM 5.3 Flash分别达到0.995和0.990的路由准确率，混合效用为0.403和0.419，后者恢复oracle的99.3\%。错路由主要发生在包含未来时态、但答案已经在文件中给出的确定问题。Qwen3-8B的路由准确率降至0.867，混合效用0.263，低于Always-Anchor的0.298。这个对比说明契约识别不是简单的时态匹配：较强基座能够利用查询与截止时点的关系，较小基座则容易把语义边界压缩为表面模式。路由实验因此支持H1在较强基座上成立，同时给出基座能力的边界。

DocFinQA和FinDVer回答第二个研究问题。DeepSeek Flash上，义务检索分支取得0.457精确匹配、平均10.69段，固定16段阅读器取得0.390、平均16段；GLM 5.3 Flash上分别为0.497对0.480，平均10.31段对16段。FinDVer上三个基座分别为0.701、0.753和0.649，均高于固定16段的0.660、0.743和0.636。四个具名RAG适配在DocFinQA上均低于义务检索分支，分支的答案覆盖率0.770和0.810也高于固定预算。规划器单独运行的成本更低，但在GLM 5.3 Flash上留下小缺口，置信度回退在24\%的问题上触发。实验支持H2：类型化义务与覆盖驱动取证在长文档上以更低证据成本取得更高准确率，并可迁移到FinDVer；收益来自证据策略，而不是阅读器或评分口径。

收入指引实验回答第三个研究问题，也给出最需要谨慎解释的结果。纯LLM锚点后验在三个基座上的准确率为0.445、0.454和0.427，决策损失为4.035、3.903和3.586，说明锚点包含非随机信号，但类别偏置和概率校准不足以直接承担低损失决策。融合后验把准确率提高到0.524、0.537和0.537，决策损失降到3.604、3.648和3.577，低于六个具名基线；在GLM 5.3 Flash和Qwen3-8B上，准确率也高于最强具名基线。纯RF诊断取得0.533、0.491、0.486和3.454，是当前最强的数值决策头。这个结果不支持“LLM锚点后验优于强RF”，但支持“语义状态与校准读出分工”：LLM负责路由、阅读和锚点状态，RF负责数值风险控制。H3应表述为融合读出相对具名基线降低决策损失，而不是LLM单独胜过RF。

消融实验进一步说明锚点层的作用与边界。去掉全部锚点后，LLM后验的宏F1从0.233、0.286、0.339降到0.213、0.211、0.272；取消硬锚点掩码后指标几乎不变，说明硬锚点在收入指引上很少触发；逐类去除锚点没有一致赢家，说明锚点类型的作用依赖基座和数据分布。输出形式对照显示，强制点答案的平均损失为4.693，固定三桶全集为4.000，损失感知门为3.454；风险收益来自集合与弃权，而不是更强的点分类器。误差分析显示，DocFinQA的错误主要来自表格块边界与算术，Qwen3-8B的错误主要来自解释性输出，AnchorEvolve的错误主要来自类别偏置，路由错误集中在语义边界。这些结果把改进方向限定在证据表示、锚点过滤和路由信号上，而不是盲目增加模型规模。

实验比较的是同基座适配，不是原系统的未修改端到端运行；标签来自XBRL核验而非人工金标准；混合负载由DocFinQA与收入指引样本构成。这些选择保证了阅读器、证据接口和评分口径可比，也意味着结论适用于给定的数据、基座和评价口径。它们不改变三个研究问题的检验结果，但界定了结果的外推范围。
